% =====================================================================
%  closure-axiomatization.tex
%
%  An equivalent axiomatization of ZF via a single Closure schema:
%  the mathematics and its machine-verified Rocq/Coq formalization.
%
%  Author: Vladimir Reshetnikov (formalization assisted by Claude Code).
%
%  - Created (UTC): 2026-06-30T17:33:23Z
%  - Repository HEAD: 9838ea000b4398278ea4b53d48ab379c0f5e6a97
%
%  Build:  lualatex closure-axiomatization.tex   (run twice for refs)
% =====================================================================
\documentclass[11pt]{article}

\usepackage{amsmath,amssymb,amsthm,mathtools}
\usepackage[margin=1in]{geometry}
\usepackage{microtype}
\usepackage{enumitem}
\usepackage{booktabs}
\usepackage{listings}
\usepackage{xcolor}
\usepackage[colorlinks=true,linkcolor=blue!55!black,urlcolor=blue!55!black,
            citecolor=blue!55!black]{hyperref}

% ---------- theorem environments ----------
\theoremstyle{plain}
\newtheorem{theorem}{Theorem}[section]
\newtheorem{lemma}[theorem]{Lemma}
\newtheorem{proposition}[theorem]{Proposition}
\newtheorem{corollary}[theorem]{Corollary}
\theoremstyle{definition}
\newtheorem{definition}[theorem]{Definition}
\newtheorem{remark}[theorem]{Remark}
\newtheorem{example}[theorem]{Example}

% ---------- notation ----------
\newcommand{\mem}{\in}
\newcommand{\Pow}{\mathcal{P}}
\newcommand{\ZF}{\mathsf{ZF}}
\renewcommand{\ZF}{\textsf{ZF}}
\newcommand{\ZFC}{\textsf{ZFC}}
\newcommand{\T}{\textsf{T}}
\newcommand{\proves}{\vdash}
\newcommand{\emptys}{\varnothing}
\DeclareMathOperator{\Sat}{Sat}
\DeclareMathOperator{\Prov}{Prov}

% ---------- Coq listing style ----------
\definecolor{coqkw}{rgb}{0.13,0.29,0.53}
\definecolor{coqcomment}{rgb}{0.40,0.40,0.40}
\definecolor{coqbg}{rgb}{0.975,0.975,0.965}

\lstdefinelanguage{Coq}{
  morekeywords={Theorem,Lemma,Definition,Fixpoint,Inductive,Hypothesis,
    Variable,Section,End,Proof,Qed,forall,exists,fun,match,with,end,
    Prop,Type,nat,list,let,in,Corollary,Check,Require,Import,From,Infix,
    Print,Assumptions},
  sensitive=true,
  morecomment=[s]{(*}{*)},
  morestring=[b]",
}
\lstset{
  language=Coq,
  basicstyle=\ttfamily\small,
  keywordstyle=\color{coqkw}\bfseries,
  commentstyle=\color{coqcomment}\itshape,
  backgroundcolor=\color{coqbg},
  frame=single,
  framerule=0.3pt,
  rulecolor=\color{black!20},
  columns=fullflexible,
  keepspaces=true,
  showstringspaces=false,
  breaklines=true,
  breakatwhitespace=true,
  inputencoding=utf8,
  extendedchars=true,
  xleftmargin=1.2em,
  xrightmargin=0.6em,
  aboveskip=0.9em,
  belowskip=0.9em,
  literate=
    {∈}{{\ensuremath{\in}}}1
    {∉}{{\ensuremath{\notin}}}1
    {⊆}{{\ensuremath{\subseteq}}}1
    {∅}{{\ensuremath{\varnothing}}}1
    {∪}{{\ensuremath{\cup}}}1
    {≺}{{\ensuremath{\prec}}}1
    {→}{{\ensuremath{\rightarrow}}}1
    {⟹}{{\ensuremath{\Longrightarrow}}}1
    {⊢}{{\ensuremath{\vdash}}}1
    {⊨}{{\ensuremath{\vDash}}}1
    {φ}{{\ensuremath{\varphi}}}1
    {ψ}{{\ensuremath{\psi}}}1
    {ω}{{\ensuremath{\omega}}}1
    {≠}{{\ensuremath{\neq}}}1
    {⋃}{{\ensuremath{\bigcup}}}1
}
% inline Coq snippet
\newcommand{\cq}[1]{\lstinline[basicstyle=\ttfamily]|#1|}

\title{\bfseries An Equivalent Axiomatization of ZF\\[2pt]
       via a Single \emph{Closure} Schema\\[6pt]
       \large Mathematics and a Machine-Verified Rocq/Coq Proof}
\author{Vladimir Reshetnikov\\[2pt]
        \small\itshape Formalization carried out with Claude Code (Rocq 9.0.1)}
\date{2026-06-30}

\begin{document}
\maketitle

\begin{abstract}
We study an alternative axiomatization $\T$ of Zermelo--Fraenkel set theory
obtained by deleting the four ``generative'' axioms --- Pairing, Union, Infinity,
and Replacement --- and adding in their place a single new schema, \emph{Closure}:
the transitive closure of any set under any \emph{set-like} definable relation is a
set. Keeping Extensionality, Regularity, Separation, and Powerset, we show
that $\T$ is \emph{exactly} $\ZF$. (Choice plays no part in the trade and is omitted
from both theories; adding it back, $\T\cup\{\mathsf C\}$ and $\ZFC$ coincide.) The
forward half (the four deleted
axioms are theorems of $\T$) is the interesting direction; it rests on one small
fact --- every set has a host, i.e.\ is a member of some set --- which turns Closure
into a fully general principle of collection. The reverse half (every $\ZF$ model
satisfies
Closure) is the textbook transitive-closure recursion. We then describe in
tutorial detail a complete machine-checked formalization in the Rocq/Coq proof
assistant, organized in four developments: a shallow embedding proving the
semantic equivalence in both directions with a ``free dependency audit'' that
certifies \emph{which} axioms each derivation actually uses; a deep embedding of
first-order logic that re-proves the forward trade from genuinely syntactic
schemas; a sound natural-deduction calculus with the cross-theory corollary
$\ZF \proves \varphi \Rightarrow \T \models \varphi$; and a from-scratch proof of
G\"odel completeness, compactness for sentence theories, and the forward
\emph{syntactic} direction $\ZF \proves \varphi \Rightarrow \T \proves \varphi$.
The whole development contains no admitted goals and rests only on the standard
axioms of classical mathematics. We close by characterizing precisely the one
direction that is true but not formalized --- the converse
$\T \proves \varphi \Rightarrow \ZF \proves \varphi$ --- and explain why it amounts
to formalizing the recursion theorem as a first-order object derivation.
\end{abstract}

\tableofcontents

% =====================================================================
\section{Introduction}\label{sec:intro}
% =====================================================================

The usual axiomatization of Zermelo--Fraenkel set theory carries several
\emph{generative} axioms whose common job is to assert that certain sets exist:
\textbf{Pairing} ($\{a,b\}$ exists), \textbf{Union} ($\bigcup s$ exists),
\textbf{Infinity} (an inductive set exists), and \textbf{Replacement} (the image
of a set under a definable function-class is a set). Each of these is, in
isolation, a recipe for collecting a previously described family of sets into one
object.

This article studies what happens when we replace all four with a \emph{single}
collection principle. Keep the ``structural'' axioms --- Extensionality,
Regularity, Separation, and Powerset --- and add the following schema.

\begin{quote}
\textbf{Closure.} For every \emph{set-like} definable binary relation $\prec$ and
every set $s$, there is a set $w$ that contains $s$ and is closed under taking
$\prec$-predecessors. (Such a $w$ is a superset of the transitive closure of $s$
under $\prec$.)
\end{quote}

Here ``set-like'' is the natural smallness condition that makes the principle
consistent: each node's class of immediate $\prec$-predecessors must already be
bounded by some set. Without it, ``$z \prec x$'' could be ``$z = z$'', and the
closure of $\{\emptyset\}$ would be the universe.

Call the resulting theory $\T$. The thesis of this article --- proved on paper and
then in full inside a proof assistant --- is:

\begin{quote}\itshape
$\T$ and $\ZF$ have exactly the same theorems.
\end{quote}

\paragraph{A word on Choice.} We work throughout in $\ZF$ rather than $\ZFC$. The
Axiom of Choice plays no part in the trade between the four generative axioms and
the Closure schema --- neither direction uses it --- so we omit it from both
theories and state everything for $\ZF$ and $\T$. Adding it back is harmless and
symmetric: since $\T$ and $\ZF$ prove the same theorems, $\T\cup\{\mathsf C\}$ and
$\ZFC$ also coincide (Corollary~\ref{cor:addC}). By the Axiom of Choice we mean any
of its standard equivalents, for instance \emph{``the Cartesian product of any
collection of nonempty sets is nonempty.''} One caveat, returned to in
\S\ref{sec:mech-shallow} and \S\ref{sec:remaining}: the \emph{proof assistant}
still uses a metatheoretic description operator (Hilbert's $\varepsilon$) to name
objects asserted to exist. That is choice in the ambient logic we reason
\emph{in}, not the set-theoretic axiom we are setting aside; the two are
independent, and we flag the metatheoretic use wherever it occurs.

What makes this more than a bookkeeping exercise is the \emph{forward} direction:
deriving Pairing, Union, Infinity, and Replacement from the lone Closure schema.
We will see that all four are instances of \emph{one} idea, and that the engine
which makes them go is not Closure by itself but Closure standing on a single,
almost-trivial principle of \textbf{hosting} --- that every set is a member of some
set:
\[
  \forall a\,\exists y\;(a \in y).
\]
(Powerset supplies it, since $a \subseteq a$ gives $a \in \Pow(a)$, but hosting is
far weaker --- see \S\ref{sec:powerset-loadbearing}.) Hosting is what lets us certify
that the relations we care about are set-like, and hence what upgrades Closure from a
statement about transitive closures into a fully general principle of collection.

\paragraph{What this article does.} There are two interleaved goals. The first is
mathematical: state the two theories precisely, prove the equivalence, exhibit the
four derivations as one schema-instance family, and isolate what genuinely powers
the trade --- the \emph{hosting} principle, a trivial shadow of Powerset --- while
showing that Powerset in full, Regularity, and Choice are all more than the trade
needs (Regularity and Choice being used in neither direction). The second is
to explain, in tutorial style, the
\emph{machine-verified} proof we built in the Rocq/Coq proof assistant. That
formalization grew well past the original equivalence: it now includes a deep
embedding of first-order logic, a sound and \emph{complete} proof calculus
(G\"odel completeness proved from scratch), a compactness theorem for sentence
theories, and the forward syntactic statement
$\ZF \proves \varphi \Rightarrow \T \proves \varphi$. We will walk through each
layer, quote the actual definitions and theorem statements, and explain the design
decisions and the one genuine obstacle that remains.

\paragraph{Roadmap.}
Sections~\ref{sec:theories}--\ref{sec:equivalence} are the mathematics: the two
theories, the forward direction (with the ``how much Powerset?'' analysis), the
reverse direction,
and the equivalence theorem. Sections~\ref{sec:mech-shallow}--\ref{sec:mech-compl}
are the formalization: the shallow embedding and the dependency audit
(\S\ref{sec:mech-shallow}), the deep embedding closing the first-order gap
(\S\ref{sec:mech-deep}), the proof calculus and soundness
(\S\ref{sec:mech-calculus}), and completeness, compactness, and the syntactic
equivalence (\S\ref{sec:mech-compl}). Section~\ref{sec:remaining} states precisely
what is proven and characterizes the single remaining direction.
Section~\ref{sec:conclusion} reflects on both punchlines.

% =====================================================================
\section{The two theories}\label{sec:theories}
% =====================================================================

We work in the first-order language with one binary relation symbol $\in$ and
equality. All schemas allow parameters (free variables other than the displayed
ones). Throughout, $\subseteq$ abbreviates the definable inclusion
$a \subseteq b \coloneqq \forall x\,(x \in a \to x \in b)$.

\subsection{The common base}

Both theories contain:

\begin{description}[leftmargin=2.2em,style=nextline]
\item[Extensionality] $\forall a\,\forall b\,\bigl(\forall x\,(x\in a \leftrightarrow x\in b)\to a=b\bigr)$.
\item[Separation (schema)] for each formula $\varphi(x,\bar p)$,
  $\forall \bar p\,\forall a\,\exists s\,\forall x\,\bigl(x\in s \leftrightarrow (x\in a \wedge \varphi(x,\bar p))\bigr)$.
\item[Powerset] $\forall a\,\exists p\,\forall x\,(x\in p \leftrightarrow x \subseteq a)$.
\item[Regularity] every nonempty set has an $\in$-minimal element.
\end{description}

(We work in $\ZF$, so Choice is not part of the common base; see the note in
\S\ref{sec:intro} and Corollary~\ref{cor:addC} for the $\ZFC$ statement.)

\subsection{The theory $\ZF$}

$\ZF$ adds the four generative axioms:

\begin{description}[leftmargin=2.2em,style=nextline]
\item[Pairing] $\forall a\,\forall b\,\exists p\,\forall x\,(x\in p \leftrightarrow (x=a \vee x=b))$.
\item[Union] $\forall s\,\exists u\,\forall x\,\bigl(x\in u \leftrightarrow \exists v\,(x\in v \wedge v\in s)\bigr)$.
\item[Infinity] there is a set containing $\emptyset$ and closed under $x \mapsto x\cup\{x\}$.
\item[Replacement (schema)] for each functional formula, the image of a set is a set.
\end{description}

\subsection{The theory $\T$}

$\T$ drops all four and adds the single \textbf{Closure} schema. Two definitions
make it precise. A binary relation $\prec$ is \emph{set-like} when
\begin{equation}\label{eq:setlike}
  \forall x\,\exists y\,\forall z\,(z \prec x \to z \in y),
\end{equation}
i.e.\ the class of immediate $\prec$-predecessors of each node is bounded by a set.
A set $w$ is \emph{closed under $\prec$-predecessors above $s$} when
\begin{equation}\label{eq:closed}
  s \subseteq w \ \wedge\ \forall u\,\forall v\,(u \prec v \wedge v \in w \to u \in w).
\end{equation}
The Closure schema asserts one such $w$ for every set-like definable $\prec$ and
every $s$:
\begin{equation}\label{eq:closure}
  \underbrace{\bigl(\forall x\,\exists y\,\forall z\,(z\prec x \to z\in y)\bigr)}_{\text{$\prec$ set-like}}
  \ \Longrightarrow\
  \forall s\,\exists w\,\underbrace{\bigl(s\subseteq w \ \wedge\ \forall u\,\forall v\,(u\prec v \wedge v\in w \to u\in w)\bigr)}_{\text{$w$ contains the transitive closure of $s$}}.
\end{equation}
There is one instance of \eqref{eq:closure} per definable binary relation $\prec$,
parameters allowed.

\begin{remark}[Why $w$ is a \emph{superset}, not the exact transitive closure]
We deliberately ask only for a set $w$ \emph{containing} the transitive closure,
not for the transitive closure itself. This is harmless: once $w$ exists,
Separation carves out the exact transitive closure as a subset of $w$. Asking only
for a superset keeps each Closure instance about an existential ``there is a
bounding set,'' which is what makes the set-like hypothesis exactly the right
amount of input.
\end{remark}

\begin{remark}[The trade is purely ``generative'']
Regularity is shared identically by both theories and (as the audit in
\S\ref{sec:mech-shallow} certifies) is used in \emph{neither} direction of the trade.
The equivalence therefore reduces to the purely ``generative'' content: over the
common base $\{\text{Ext},\text{Sep},\text{Pow}\}$ (with Regularity along for the
ride),
the four axioms $\{\text{Pairing},\text{Union},\text{Infinity},\text{Replacement}\}$
are interderivable with the single schema Closure.
\end{remark}

% =====================================================================
\section{The forward direction: \texorpdfstring{$\T \proves \ZF$}{T proves ZF}}\label{sec:forward}
% =====================================================================

This is the interesting half. We assume only
$\{\text{Extensionality},\text{Separation},\text{Powerset},\text{Closure}\}$ and
derive Pairing, Union, Replacement, and Infinity. (Regularity is not
needed here; it is a genuine passenger, as the dependency audit in
\S\ref{sec:mech-shallow} will certify.)

\subsection{The linchpin: every set has a host}\label{sec:linchpin}

The engine of the forward direction is one almost-trivial fact: every set is a
\emph{member} of some set. Call a set $y$ a \emph{host} of $a$ when $a \in y$. As we
will see, the trade needs nothing about $a$ beyond its \emph{having} a host.

Powerset is one ready source of hosts:

\begin{lemma}[self-membership in the powerset]\label{lem:self}
For every set $a$, $\ a \in \Pow(a)$.
\end{lemma}
\begin{proof}
$a \subseteq a$ holds trivially, and by Powerset $x \in \Pow(a) \leftrightarrow x \subseteq a$;
take $x = a$.
\end{proof}

It is worth seeing how little of Powerset this uses. Consider any class relation
$\prec$ that is \emph{singleton-valued} above each node --- i.e.\ for each $x$ there
is at most one $z$ with $z \prec x$, given by some class function $z = G(x)$. Is such
a $\prec$ set-like? The set-like condition \eqref{eq:setlike} asks for a set $y$
bounding $\{\,z : z \prec x\,\} = \{G(x)\}$ --- that is, a \emph{host of $G(x)$}.
Lemma~\ref{lem:self} supplies one ($\Pow(G(x))$), but any host would do; the
powerset is never opened up. More generally, whenever the predecessor-class above
each node already sits inside some set, that set is the bound.

\begin{quote}\itshape
This is the conceptual heart of the whole equivalence: \textbf{hosting --- ``every
set is a member of some set'' --- turns ``has a definable predecessor'' into ``is
set-like.''} Closure then collects the predecessors. The four generative axioms are
four ways of choosing the relation $\prec$ and the seed $s$. We will see in
\S\ref{sec:powerset-loadbearing} that hosting is \emph{all} that is used, so full
Powerset is more than the trade needs.
\end{quote}

\subsection{Four derivations, one schema-instance family}

\begin{center}
\renewcommand{\arraystretch}{1.35}
\begin{tabular}{@{}lll@{}}
\toprule
Axiom & Relation $\prec$ (read $z \prec x$) & Seed $s$ \\
\midrule
Union       & $z \in x$ \hfill (membership) & the given set $s$ \\
Replacement & $x \in a \wedge z = F(x)$ \hfill (function graph) & the domain $a$ \\
Pairing     & $(x{=}\emptyset \wedge z{=}a)\vee(x{=}\{\emptyset\}\wedge z{=}b)$ & the 2-node seed $\{\emptyset,\{\emptyset\}\}$ \\
Infinity    & $z = x \cup \{x\}$ \hfill (successor) & $\{\emptyset\}$ \\
\bottomrule
\end{tabular}
\end{center}

In each case the recipe is the same: (i) check the relation is set-like (using
Lemma~\ref{lem:self}); (ii) apply Closure to get a closed superset $w$; (iii) use
Separation to carve the exact set we wanted out of $w$.

\paragraph{Union} is closure under $\in$. The relation $z \prec x \coloneqq z \in x$
is set-like because the predecessors of $x$ are exactly the elements of $x$,
bounded by $x$ itself ($y = x$). Closure applied to $s$ gives a $w$ with
$s \subseteq w$ and $w$ closed under $\in$-predecessors. Hence every element of
every member of $s$ lies in $w$, and Separation extracts
$\{\,x \in w : \exists v\,(x \in v \wedge v \in s)\,\} = \bigcup s$.

\paragraph{Replacement} is closure along a function graph. For a class function
$F$ and domain $a$, take $z \prec x \coloneqq x \in a \wedge z = F(x)$. The
predecessors of $x$ form the singleton $\{F(x)\}$ (or $\emptyset$ if
$x \notin a$), bounded by any host of $F(x)$ (Lemma~\ref{lem:self} gives one,
$\Pow(F(x))$) --- this is where hosting is used. Closure applied to $a$ yields a $w$
containing the whole image
$\{F(x) : x \in a\}$, and Separation carves out exactly that image.

\paragraph{Pairing} is the subtle one, and it is precisely where the host-providing
role becomes visible as a \emph{design constraint}. We want
$\{a,b\}$ for arbitrary $a,b$. The naive attempt --- one seed node whose
predecessors are $\{a,b\}$ --- is self-defeating: to certify that relation
set-like we would need a set bounding $\{a,b\}$, which is Pairing itself. The fix
is to \emph{split $a$ and $b$ across two different seed nodes}. Use the two-node
seed $\{\emptyset,\{\emptyset\}\}$ (a fixed two-element set, built without Pairing)
and the relation
\[
  z \prec x \ \coloneqq\ (x = \emptyset \wedge z = a)\ \vee\ (x = \{\emptyset\} \wedge z = b).
\]
Now each node has a \emph{singleton} predecessor-class --- $\{a\}$ above
$\emptyset$, $\{b\}$ above $\{\emptyset\}$ --- so hosting makes it set-like (any
hosts of $a$ and $b$ work --- e.g.\ $\Pow(a)$, $\Pow(b)$), with no circular appeal
to Pairing. Closure of the seed then drags both $a$ and $b$ into one set $w$, and
Separation gives $\{x \in w : x = a \vee x = b\} = \{a,b\}$.

\paragraph{Infinity} is closure of $\{\emptyset\}$ under the successor relation
$z \prec x \coloneqq z = x \cup \{x\}$. The successor is unique, so the
predecessor-class is again a singleton, hostable (any host of the successor works);
the relation is set-like. Closure applied to $\{\emptyset\}$ gives a set $w$ containing
$\emptyset$ and closed under successors --- an inductive set, which is exactly what
Infinity asserts. (Once Pairing and Union are available, the successor
$x \cup \{x\}$ is itself a set, so the relation is genuinely definable.)

\subsection{How much Powerset? Only hosting}\label{sec:powerset-loadbearing}

The four derivations use Powerset in a single, uniform role: as a source of hosts in
Lemma~\ref{lem:self}, the fact that makes singleton-valued (and more generally
suitably bounded) relations set-like. Union is the lone exception --- its relation
$z \in x$ is bounded by $x$ directly, with no host needed --- which is why the audit
in \S\ref{sec:mech-shallow} shows Union depending on Separation and Closure
\emph{only}.

This answers a natural question: \emph{do we need the full Powerset axiom?} No. The
trade never opens a powerset up; it uses only the consequence
\[
  (\mathrm H)\qquad \forall a\,\exists y\;(a \in y)
  \qquad\text{(``every set is a member of some set''),}
\]
the \emph{hosting} principle. Powerset yields $(\mathrm H)$ via $a \in \Pow(a)$, but
$(\mathrm H)$ is far weaker --- it holds, for instance, in $H(\aleph_1)$, the
hereditarily countable sets, which has no $\Pow(\omega)$. The machine confirms it:
with the hypothesis weakened from Powerset to $(\mathrm H)$, all four derivations
still go through, and \cq{Check} reports them depending on $(\mathrm H)$, not on
Powerset (\S\ref{sec:mech-shallow}).

\begin{remark}[Would \emph{finite} Powerset suffice? No]\label{rem:finitepow}
A tempting weakening is ``Powerset restricted to finite sets'' (under any reasonable
notion of finite --- Dedekind, Tarski). It is not enough. The predecessors we must
host --- $a, b$ in Pairing, $F(x)$ in Replacement, $x \cup \{x\}$ in Infinity --- are
\emph{arbitrary}, possibly infinite sets, and finite Powerset produces $\Pow(a)$ only
for finite $a$, so it cannot host an infinite set; it does not yield $(\mathrm H)$.
Hosting infinite sets is moreover unavoidable: deriving Pairing of two infinite sets
forces, through set-likeness, a bounding set containing each of them, and the base
$\{$Extensionality, Separation, Closure$\}$ has no other way to manufacture a host
(Closure cannot host a set it is not already handed, on pain of circularity). Indeed
Pairing already implies $(\mathrm H)$ --- $a \in \{a,b\}$ --- so any sufficient
fragment must prove $(\mathrm H)$, which finite Powerset does not. The right way to
weaken Powerset here is therefore not ``on finite sets'' but ``replace it by
hosting.''
\end{remark}

So, on the structural side, the picture is uniform: \emph{none} of Powerset,
Regularity, or Choice is used in full. Regularity and Choice are passengers, used in
neither direction (\S\ref{sec:reverse}); Powerset is used only through its trivial
hosting shadow $(\mathrm H)$, and only forward (the reverse needs no Powerset at
all). The genuine engine of the equivalence is just
$\{$Extensionality, Separation, Closure$\}$ together with hosting.

% =====================================================================
\section{The reverse direction: \texorpdfstring{$\ZF \proves \text{Closure}$}{ZF proves Closure}}\label{sec:reverse}
% =====================================================================

The reverse direction is the standard transitive-closure construction; the only
care needed is in collecting an $\omega$-indexed family of sets into one object,
which is exactly the job Replacement and Infinity exist to do.

Fix a set-like relation $R$ and a set $s$. Define the one-step operator
\[
  g(t) \ =\ t \ \cup\ \{\,u : \exists v \in t,\ u \mathrel{R} v\,\}.
\]
The bracketed predecessor-set is genuinely a set. Set-likeness says that for each
$v$ there \emph{exists} a set bounding $\{u : u \mathrel R v\}$; the \emph{Collection}
schema --- a theorem of $\ZF$ (provable from Replacement together with Foundation,
and itself choice-free) --- gathers these into a single set $B$ holding a bound for
every $v \in t$, so $\bigcup B$ bounds all predecessors at once and Separation carves
out the exact predecessor-set. (Collection delivers a bounding \emph{set}, not a
chosen function, so no Choice is needed --- cf.\ the note in \S\ref{sec:intro}.) So
$g : V \to V$ is a definable total operation.

Now iterate: $W_0 = s$ and $W_{n+1} = g(W_n)$, indexing on the \emph{metatheory's}
natural numbers. The closure we want is $w = \bigcup_n W_n$. To produce $w$ as a
single set we must turn the meta-indexed family $\{W_n\}$ into an
\emph{object}-indexed family. Use Infinity to obtain an inductive set, and inside it
the von~Neumann numerals $\underline n$ ($\underline 0 = \emptyset$,
$\underline{n+1} = \underline n \cup \{\underline n\}$). The assignment
$\underline n \mapsto W_n$ is a definable class function on the numerals, so
Replacement makes its image $\{W_n : n\}$ a set, and Union gives
$w = \bigcup_n W_n$.

For this to be well defined we need the numerals to be \emph{distinct}, i.e.\
$\underline m = \underline n \Rightarrow m = n$, so that ``$\underline n \mapsto
W_n$'' is unambiguous. Distinctness follows from $\underline m \in \underline n$
whenever $m < n$, together with the fact that no \emph{numeral} is a member of
itself. One might reach for Regularity here (apply it to $\{\underline n\}$), but
that is a \emph{convenience}, not a necessity. The finite von~Neumann numerals are
non-self-membered already in $\ZF$ \emph{without} Foundation: each $\underline n$ is
a genuine finite ordinal --- a transitive set strictly well-ordered by $\in$ --- and
an ordinal is never a member of itself (that would violate the irreflexivity of its
own $\in$-ordering), a fact established by a plain induction on $n$ rather than by
the global Foundation axiom. So Regularity is \emph{not} needed for the reverse
direction; we make this precise, and let the machine certify it, in
\S\ref{sec:mech-shallow}. Finally, $w$ contains $s = W_0$ and is closed under
$R$-predecessors: if $u \mathrel R v$ and $v \in w$, then $v \in W_n$ for some
$n$, hence $u \in g(W_n) = W_{n+1} \subseteq w$.

\begin{remark}[Regularity is a passenger, in both directions]
\label{rem:mirror}
What actually powers the reverse construction is Replacement, Union, and Infinity (to
collect the family) --- not Powerset, and (by the Foundation-free argument above) not
Regularity either. The machine audit in \S\ref{sec:mech-shallow} confirms both
omissions, yielding the sharper statement
\[
  \ZF - \text{Powerset} - \text{Regularity} \ \proves\ \text{Closure}.
\]
Together with the forward audit (which uses neither Regularity nor Choice), this says
that \textbf{Regularity, like Choice, is a passenger of the whole trade}: the
interderivability of Closure with $\{$Pairing, Union, Infinity, Replacement$\}$ holds
over the bare base $\{$Extensionality, Separation, Powerset$\}$, with Regularity and
Choice added equally to both sides and used by neither. The one axiom that is
genuinely \emph{asymmetric} is Powerset --- and even it is used forward only through
its hosting shadow (\S\ref{sec:powerset-loadbearing}) and is idle in reverse.
\end{remark}

% =====================================================================
\section{The equivalence theorem}\label{sec:equivalence}
% =====================================================================

\begin{theorem}[Equivalence]\label{thm:equiv}
Over the base $\{\text{Extensionality}, \text{Separation}, \text{Powerset}\}$, the
schema Closure is interderivable with $\{\text{Pairing}, \text{Union},
\text{Infinity}, \text{Replacement}\}$. Consequently --- adding Regularity to both
sides, where it is used by neither --- the theory $\T$ (adding Closure) and the
theory $\ZF$ (adding Pairing, Union, Infinity, Replacement) have exactly the same
theorems.
\end{theorem}
\begin{proof}
$\T \proves \ZF$: Pairing, Union, Replacement, and Infinity are
Theorems by the four derivations of \S\ref{sec:forward}, none of which uses any
$\ZF$-only axiom. $\ZF \proves \T$: Closure is a Theorem by the construction of
\S\ref{sec:reverse}. Since the two theories share their remaining axioms verbatim,
they prove the same sentences.
\end{proof}

\begin{corollary}[Adding Choice]\label{cor:addC}
$\T \cup \{\mathsf C\}$ and $\ZFC$ have exactly the same theorems.
\end{corollary}
\begin{proof}
Choice is the \emph{same} sentence $\mathsf C$ adjoined to both theories. For any
$\varphi$, the deduction theorem gives $\T \cup \{\mathsf C\} \proves \varphi$ iff
$\T \proves (\mathsf C \to \varphi)$ iff (Theorem~\ref{thm:equiv})
$\ZF \proves (\mathsf C \to \varphi)$ iff $\ZFC \proves \varphi$. The argument is
insensitive to which standard form of $\mathsf C$ one adopts.
\end{proof}

The rest of the article is about making Theorem~\ref{thm:equiv} --- and a good deal
more --- checkable by machine.

% =====================================================================
\section{Mechanization I: the shallow embedding and the dependency audit}\label{sec:mech-shallow}
% =====================================================================

The formalization lives in the \texttt{src/SetTheory/} directory of the
repository, in four Rocq/Coq source files. This section covers the two that prove
the semantic equivalence directly: \texttt{Forward.v} and \texttt{Reverse.v}.

\subsection{Why Rocq/Coq, and what a ``shallow embedding'' is}

We use Rocq/Coq (version 9.0.1). The decisive feature is the
\textbf{Section/Hypothesis} idiom, which fits our problem like a glove. A Coq
\cq{Section} lets us fix an abstract structure and assume axioms about it as local
\cq{Hypothesis}es; every theorem proved inside is, after the section closes,
automatically generalized over the structure \emph{and over exactly the hypotheses
it used}. That is precisely the shape of a metatheorem of the form ``every model of
these axioms satisfies that conclusion.''

The structure is modelled \emph{shallowly}: the universe of sets is an abstract Coq
type \cq{V}, membership is an abstract relation \cq{mem : V -> V -> Prop}, and the
axiom schemas are rendered with the metatheory's own predicates. For instance
Separation quantifies over a Coq predicate \cq{P : V -> Prop}:

\begin{lstlisting}
Variable V : Type.
Variable mem : V -> V -> Prop.
Infix "∈" := mem (at level 70, no associativity).
Definition Sub (a b : V) : Prop := forall x, x ∈ a -> x ∈ b.

Hypothesis Extensionality :
  forall a b, (forall x, x ∈ a <-> x ∈ b) -> a = b.
Hypothesis Separation :
  forall (a : V) (P : V -> Prop),
    exists s, forall x, x ∈ s <-> (x ∈ a /\ P x).
Hypothesis Powerset :
  forall a, exists p, forall x, x ∈ p <-> Sub x a.
\end{lstlisting}

The Closure schema becomes a quantification over a Coq relation
\cq{R : V -> V -> Prop}, with set-likeness defined exactly as in
\eqref{eq:setlike}:

\begin{lstlisting}
Definition SetLike (R : V -> V -> Prop) : Prop :=
  forall x, exists y, forall z, R z x -> z ∈ y.

Hypothesis Closure :
  forall R : V -> V -> Prop, SetLike R ->
    forall s, exists w, Sub s w /\ (forall u v, R u v -> v ∈ w -> u ∈ w).
\end{lstlisting}

\paragraph{What a shallow embedding proves, and its faithfulness.} Because the
schemas range over \emph{all} Coq predicates/relations, the artifact establishes
the \emph{second-order} statement: every structure $(V,\in)$ satisfying these
axioms also satisfies the derived ones. This is the standard, accepted way to
mechanize an equivalence of axiomatizations, and it is strictly stronger than the
first-order schematic reading on the model side. It stays faithful to the genuine
first-order claim because \emph{every derivation instantiates a schema at one
concrete, definable relation} --- the membership relation, a function graph, the
two-branch pairing relation, the successor relation --- exactly the instances a
first-order proof would write down. (Section~\ref{sec:mech-deep} removes even this
caveat for the forward direction, by re-proving it with schemas that range over
syntactic formulas.)

\subsection{The linchpin and the four theorems, in Coq}

The linchpin is two lines of Coq --- the \cq{host} operator (from the \cq{Hosting}
hypothesis $\forall a\,\exists y\,a\in y$, packaged by description) and its spec ---
exactly the pivot the informal argument advertised:

\begin{lstlisting}
(* THE LINCHPIN, minimal form: every set has a host (a ∈ host a). *)
Definition host (a : V) : V :=
  proj1_sig (constructive_indefinite_description _ (Hosting a)).
Lemma host_spec : forall a, a ∈ host a.
Proof. intro a. exact (proj2_sig (constructive_indefinite_description _ (Hosting a))). Qed.
\end{lstlisting}

A side lemma \cq{powerset_gives_hosting} records that Powerset satisfies
\cq{Hosting} ($a \in \Pow(a)$) --- the only place the file touches the Powerset
hypothesis at all.

The four relations of the schema-instance family are declared verbatim:

\begin{lstlisting}
Definition pairRel (a b : V) : V -> V -> Prop :=
  fun z x => (x = emptyset /\ z = a) \/ (x = single_empty /\ z = b).
Definition memRel : V -> V -> Prop := fun z x => z ∈ x.
Definition graphRel (F : V -> V) (a : V) : V -> V -> Prop :=
  fun z x => x ∈ a /\ z = F x.
Definition succRel : V -> V -> Prop :=
  fun z x => forall t, t ∈ z <-> (t ∈ x \/ t = x).
\end{lstlisting}

Each of Pairing, Union, Replacement, Infinity is then proved by the three-step
recipe (set-like check $\to$ Closure $\to$ Separation). Union is the shortest and
shows the pattern cleanly:

\begin{lstlisting}
Theorem Union :
  forall s, exists u, forall x, x ∈ u <-> exists v, x ∈ v /\ v ∈ s.
Proof.
  intro s.
  assert (HSL : SetLike memRel).
  { intro x. exists x. intros z Hz. unfold memRel in Hz. exact Hz. }
  destruct (Closure memRel HSL s) as [w [Hsub Hclosed]].
  exists (sep w (fun x => exists v, x ∈ v /\ v ∈ s)).
  intro x. split.
  - intro H. exact (sep_elim2 _ _ _ H).
  - intro H. apply sep_intro.
    + destruct H as [v [Hxv Hvs]]. apply (Hclosed x v).
      * unfold memRel. exact Hxv.
      * apply Hsub. exact Hvs.
    + exact H.
Qed.
\end{lstlisting}

Read it as English: ``\cq{memRel} is set-like (host $= x$ itself); apply
\cq{Closure} to get a closed superset $w$ of $s$; separate out the genuine union
elements; membership in either direction is immediate from
closure/Separation.''

The Pairing proof is where the two-node seed appears and where the set-like check
does real work: it splits on whether the node is $\emptyset$ or $\{\emptyset\}$ and
supplies \cq{host a} or \cq{host b} accordingly --- the Coq transcription of the
``split $a$ and $b$ across different nodes'' argument from \S\ref{sec:forward}. (The
seed $\{\emptyset,\{\emptyset\}\}$ is itself built from \cq{host} and a one-edge
Closure, not from Powerset.)

\subsection{The free dependency audit}

Here is the part that turns the ``only hosting is used'' claim from rhetoric into a
machine certificate. When a Coq \cq{Section} closes, each theorem is generalized over
precisely the hypotheses it referenced; the trailing \cq{Check} commands print those
generalized types, so we can simply \emph{read off} which axioms each derivation
used:

\begin{lstlisting}
End ClosureEquivalence_Forward.
Check Pairing.
Check Union.
Check Replacement.
Check Infinity.
\end{lstlisting}

The printed signatures certify:

\begin{itemize}[leftmargin=1.4em]
\item \textbf{Union} depends on \emph{Separation and Closure only} --- not
  Hosting, not Powerset, not Extensionality, not the nonempty-domain assumption.
  (Its relation $z\in x$ is self-bounding.)
\item \textbf{Replacement} depends on \emph{Separation, Hosting, and Closure} ---
  Hosting in exactly its host-providing role, and \emph{not} Powerset (nor
  Extensionality, nor nonemptiness).
\item \textbf{Pairing} depends on \emph{Separation, Hosting, Closure} and a
  \emph{nonempty domain} (to build $\emptyset$) --- and, now that the seed is built
  by Closure rather than by a powerset, not even Extensionality.
\item \textbf{Infinity} additionally needs \emph{Extensionality} (for the
  inductive-set uniqueness step).
\item \textbf{Powerset} is referenced by \emph{none} of the four (only by the side
  lemma \cq{powerset_gives_hosting}); \textbf{Regularity} by none either --- both,
  with Choice, are genuine passengers.
\end{itemize}

The same audit on \texttt{Reverse.v} (\cq{Check Closure_holds}) shows
\cq{Closure_holds} depending on a nonempty domain, Extensionality, Separation,
Pairing, Union, Infinity, and Replacement --- but \textbf{neither Powerset nor
Regularity}. This is the machine confirming the sharper
$\ZF - \text{Powerset} - \text{Regularity} \proves \text{Closure}$ of
Remark~\ref{rem:mirror}, and with it the asymmetry: Powerset is used forward (only
through hosting, \S\ref{sec:powerset-loadbearing}) and idle in reverse, while
Regularity is idle in \emph{both} directions --- a passenger of the trade alongside
Choice.

\subsection{The reverse construction in Coq}

\texttt{Reverse.v} transcribes \S\ref{sec:reverse} directly. The meta-level
iteration is an ordinary Coq \cq{Fixpoint} on \cq{nat}:

\begin{lstlisting}
Fixpoint iterate (g : V -> V) (s : V) (n : nat) : V :=
  match n with O => s | S k => g (iterate g s k) end.
\end{lstlisting}

The numerals' distinctness is established \emph{without} Regularity: each numeral is
a transitive set (\cq{onat_trans}), hence non-self-membered (\cq{onat_no_self}), both
by induction on \cq{nat}, from which injectivity follows:

\begin{lstlisting}
Lemma onat_trans   : forall n x, x ∈ onat n -> Sub x (onat n).   (* Foundation-free *)
Lemma onat_no_self : forall n, ~ onat n ∈ onat n.                (* no Regularity *)
Lemma onat_inj     : forall i j, onat i = onat j -> i = j.
\end{lstlisting}

\cq{onat_inj} is what makes the map ``numeral $\mapsto$ iterate'' well defined, so
Replacement can collect $\{W_n\}$ into a set and Union flatten it. Because this path
no longer touches Regularity, \cq{Check Closure_holds} drops it from the dependency
list (Remark~\ref{rem:mirror}). The headline theorem is then

\begin{lstlisting}
Theorem Closure_holds :
  forall R : V -> V -> Prop, SetLike R ->
    forall s, exists w, Sub s w /\ (forall u v, R u v -> v ∈ w -> u ∈ w).
\end{lstlisting}

with $w = \bigcup_n W_n$ realized as \cq{ounion (imageR (Ffun R HSL s) Inf)} ---
the object union of the Replacement-image, over the inductive set \cq{Inf}, of the
numeral-to-iterate map \cq{Ffun}.

\subsection{Where classical logic --- and metatheoretic choice --- enter}

Both files import \cq{ClassicalEpsilon}. It is used to \emph{package} the
existential axioms as operations --- turning ``$\exists p,\ \dots$'' (Powerset)
into a function \cq{power : V -> V} with a specification lemma, via
\cq{constructive_indefinite_description} --- and, in \texttt{Reverse.v}, to extract
a bounding function \cq{boundf} from set-likeness and to index the iteration.

It is worth being precise about what kind of choice this is, since we went to some
trouble in \S\ref{sec:theories} to keep the \emph{set-theoretic} Axiom of Choice out
of both theories. The description operator here is choice in the \emph{ambient
logic} we reason \emph{in} --- it only gives a name to an object the theory already
asserts to exist --- and it is independent of the object-level Axiom of Choice,
present in both directions regardless. In particular \cq{boundf} is a
\emph{convenience}: as \S\ref{sec:reverse} explained, the underlying $\ZF$
mathematics needs only the choice-free Collection schema to form the
predecessor-set, so no set-theoretic Choice is smuggled in through the back door.
The \cq{Print Assumptions} audit in \S\ref{sec:remaining} confirms that neither
direction depends on an object-level Choice axiom.

% =====================================================================
\section{Mechanization II: closing the first-order gap}\label{sec:mech-deep}
% =====================================================================

The shallow embedding's one honest caveat is that its schemas range over arbitrary
Coq predicates rather than over first-order \emph{formulas}. \texttt{Deep.v}
removes that caveat for the forward direction by building a genuine deep embedding
of the first-order language of set theory and re-proving the trade from schemas
that quantify over syntactic formulas.

\subsection{Syntax, environments, and Tarski satisfaction}

A formula is an element of an inductive type, with De~Bruijn indices for variables
(so a variable is a \cq{nat} counting binders outward, and $\alpha$-equivalence is
syntactic equality):

\begin{lstlisting}
Inductive form : Type :=
| fMem : nat -> nat -> form
| fEq  : nat -> nat -> form
| fBot : form
| fImp : form -> form -> form
| fAnd : form -> form -> form
| fOr  : form -> form -> form
| fAll : form -> form
| fEx  : form -> form.
\end{lstlisting}

An \emph{environment} \cq{e : nat -> V} assigns a domain element to each variable.
\cq{scons d e} pushes a new binding \cq{d} at index $0$ and shifts the rest ---
this is how a quantifier extends the environment under its binder. Tarski
satisfaction is then a straightforward structural recursion:

\begin{lstlisting}
Definition scons (d : V) (e : nat -> V) : nat -> V :=
  fun n => match n with O => d | S k => e k end.

Fixpoint Sat (e : nat -> V) (f : form) {struct f} : Prop :=
  match f with
  | fMem i j => (e i) ∈ (e j)
  | fEq i j  => e i = e j
  | fBot     => False
  | fImp a b => Sat e a -> Sat e b
  | fAnd a b => Sat e a /\ Sat e b
  | fOr a b  => Sat e a \/ Sat e b
  | fAll a   => forall d, Sat (scons d e) a
  | fEx a    => exists d, Sat (scons d e) a
  end.
\end{lstlisting}

Two infrastructural lemmas carry the whole development. \cq{Sat_ext} says
satisfaction depends only on the environment's values on the variables that occur;
\cq{Sat_rename} relates satisfaction of a renamed formula to satisfaction under the
renamed environment:
\[
  \Sat\ e\ (\mathtt{rename}\ r\ f)\ \leftrightarrow\ \Sat\ (e\circ r)\ f.
\]
Renaming is the deep-embedding analogue of variable substitution; because our
signature is \emph{purely relational}, there are no function symbols and hence no
compound terms --- a ``term'' is just a variable --- so substitution of a variable
for a variable is exactly a renaming. This single observation keeps the entire
metatheory free of a general substitution operation, which is usually the most
painful part of a deep embedding.

\subsection{First-order Separation and Closure}

Now the schemas range over formulas. A formula \cq{phi} with free variable $0$
denotes a unary predicate (its other free variables are parameters, read from the
environment \cq{e}); a formula \cq{psi} with free variables $0,1$ denotes a binary
relation \cq{relOf psi e}:

\begin{lstlisting}
Hypothesis SeparationFO :
  forall (phi : form) (e : nat -> V) (a : V),
    exists s, forall x, x ∈ s <-> (x ∈ a /\ Sat (scons x e) phi).

Definition relOf (psi : form) (e : nat -> V) : V -> V -> Prop :=
  fun z x => Sat (scons z (scons x e)) psi.

Hypothesis ClosureFO :
  forall (psi : form) (e : nat -> V),
    SetLike (relOf psi e) ->
    forall s, exists w, Sub s w /\ (forall u v, relOf psi e u v -> v ∈ w -> u ∈ w).
\end{lstlisting}

\texttt{Deep.v} then re-derives Pairing, Union, first-order Replacement, and
Infinity from \emph{these} schemas. The new obligation, compared to the shallow
version, is that each relation used must be exhibited as a concrete \cq{form} and
its \cq{Sat}-meaning verified. For the simple cases the defining formula is short;
for Replacement, expressing ``$\exists x \in a,\ \psi(y,x)$'' as a formula requires
shuffling De~Bruijn indices, which is exactly what \cq{Sat_rename} is for. The
upshot is a \emph{formal} certificate --- not a meta-remark --- that the
schema-trade uses only first-order definable instances:

\begin{lstlisting}
Theorem Pairing       : forall a b, exists p, forall x, x ∈ p <-> (x = a \/ x = b).
Theorem Union         : forall s, exists u, forall x, x ∈ u <-> exists v, x ∈ v /\ v ∈ s.
Theorem ReplacementFO : (* image of a set under a formula-defined functional relation *)
Theorem Infinity      : (* an inductive set *)
\end{lstlisting}

\paragraph{Why only the forward direction deepens.} The forward trade ports to the
deep setting cheaply because every relation it uses is a short first-order formula.
The reverse direction does \emph{not}: its essential use of Replacement collects
the family $\{W_n\}$ via the map $\underline n \mapsto W_n$, and that map is
first-order definable only through the \emph{recursion theorem}. \texttt{Reverse.v}
deliberately sidesteps the recursion theorem by iterating on the meta-level
\cq{nat}; rendering the same construction first-order would mean formalizing the
recursion theorem's defining formula --- a heavy object-level construction. So the
genuine first-order content of the reverse direction \emph{is} the recursion
theorem, and we leave it at the (already verified) shallow level. We return to this
point in \S\ref{sec:remaining}.

% =====================================================================
\section{Mechanization III: a proof calculus, soundness, and a cross-theory bridge}\label{sec:mech-calculus}
% =====================================================================

So far everything has been semantic (about models). \texttt{Deep.v} also builds a
\emph{syntactic} layer: an actual proof calculus over \cq{form}, a soundness
theorem, and a first bridge between the two theories at the level of provability.

\subsection{A natural-deduction calculus}

\cq{Prov G a} means ``$a$ is derivable from the context \cq{G} (a list of
formulas).'' The calculus is classical natural deduction:

\begin{lstlisting}
Inductive Prov : list form -> form -> Prop :=
| P_ass    : forall G a, In a G -> Prov G a
| P_impI   : forall G a b, Prov (a :: G) b -> Prov G (fImp a b)
| P_impE   : forall G a b, Prov G (fImp a b) -> Prov G a -> Prov G b
| P_botE   : forall G a, Prov G fBot -> Prov G a
| P_lem    : forall G a, Prov G (fOr a (fImp a fBot))
| P_andI   : forall G a b, Prov G a -> Prov G b -> Prov G (fAnd a b)
| P_andE1  : forall G a b, Prov G (fAnd a b) -> Prov G a
| P_andE2  : forall G a b, Prov G (fAnd a b) -> Prov G b
| P_orI1   : forall G a b, Prov G a -> Prov G (fOr a b)
| P_orI2   : forall G a b, Prov G b -> Prov G (fOr a b)
| P_orE    : forall G a b c, Prov G (fOr a b) -> Prov (a :: G) c -> Prov (b :: G) c -> Prov G c
| P_allI   : forall G a, Prov (map (rename S) G) a -> Prov G (fAll a)
| P_allE   : forall G a k, Prov G (fAll a) -> Prov G (rename (inst k) a)
| P_exI    : forall G a k, Prov G (rename (inst k) a) -> Prov G (fEx a)
| P_exE    : forall G a c, Prov G (fEx a) -> Prov (a :: map (rename S) G) (rename S c) -> Prov G c
| P_eqRefl : forall G k, Prov G (fEq k k)
| P_eqElim : forall G i j a,
    Prov G (fEq i j) -> Prov G (rename (inst i) a) -> Prov G (rename (inst j) a).
\end{lstlisting}

The De~Bruijn discipline shows up in the quantifier rules. In \cq{P_allI}
(universal introduction), to prove $\forall a$ we shift the context outward by
\cq{rename S} --- index $0$ becomes the fresh ``eigenvariable'' bound by the new
$\forall$, and the old context variables move up by one. In \cq{P_allE} (universal
elimination), instantiating $\forall a$ at variable $k$ is the renaming
\cq{rename (inst k) a}, where \cq{inst k} maps the bound index $0$ to $k$ and
decrements the rest. The existential rules are dual. Crucially, instantiation is
\emph{just renaming} --- the purely relational signature again paying off ---
so the same \cq{rename}/\cq{Sat_rename} machinery handles quantifiers with no
separate substitution.

\subsection{The equality rule: a genuine correctness fix}

The two equality rules deserve a word, because getting them right was a real
correction during development. The reflexivity rule \cq{P_eqRefl} is obvious. For
the elimination rule, a tempting but \emph{too weak} design is a
``replace-the-variable-everywhere'' congruence: from $i = j$ and $\varphi$, infer
$\varphi$ with every occurrence of $i$ rewritten to $j$. That rule cannot even
derive \emph{symmetry} of equality: to get $j = i$ from $i = j$ you would rewrite
$i$ to $j$ in the formula $j = i$, but $j = i$ contains no $i$ in the right slot to
key on, and the rewrite erases the very occurrence you need.

The fix is the proper \textbf{Leibniz} rule \cq{P_eqElim}: a formula \cq{a} is
treated as a property with a hole at De~Bruijn index $0$; from $i = j$ and
``\cq{a} with the hole filled by $i$'' (\cq{rename (inst i) a}) we infer ``\cq{a}
with the hole filled by $j$'' (\cq{rename (inst j) a}). This single rule yields
symmetry, transitivity, and full congruence, and --- as we verify --- it is sound.
Without it the completeness theorem of \S\ref{sec:mech-compl} would have been
unreachable, since the term model is a quotient by provable equality and needs all
the equality reasoning. A small example of the rule earning its keep: it is exactly
what makes the term-model congruence lemmas go through.

\subsection{Soundness, and the bridge $\ZF \proves \varphi \Rightarrow \T \models \varphi$}

Soundness says the calculus only proves things that are true in every model
satisfying the context:

\begin{lstlisting}
Theorem soundness :
  forall G a, Prov G a -> forall e, (forall x, In x G -> Sat e x) -> Sat e a.
\end{lstlisting}

The proof is one case per rule; the quantifier and equality cases are where
\cq{Sat_rename} and \cq{Sat_ext} are spent.

Now encode the $\ZF$ axioms as closed formulas (\cq{Ext_form}, \cq{Pair_form},
\dots, \cq{Reg_form}, with Separation and Replacement as schemas over \cq{form}),
collect them into a predicate \cq{ZFax}, and define provability from the $\ZF$
axioms:

\begin{lstlisting}
Definition ZFprov (phi : form) : Prop :=
  exists G, (forall x, In x G -> ZFax x) /\ Prov G phi.

Theorem ZF_provable_holds_in_T :
  forall phi, ZFprov phi -> forall e, Sat e phi.
\end{lstlisting}

Because the ambient section assumes the $\T$-axioms, this theorem says: \emph{in
any model of the Closure axiomatization $\T$, every $\ZF$-provable formula holds},
i.e.\ $\ZF \proves \varphi \Rightarrow \T \models \varphi$. It is the marriage of
syntactic soundness with the forward semantic equivalence: each encoded $\ZF$ axiom
is checked to hold in the $\T$-model (using the derived \cq{Pairing}, \cq{Union},
\cq{ReplacementFO}, \cq{Infinity} for the four traded ones, and the $\T$-hypotheses
for the shared ones), and then \cq{soundness} propagates truth through any $\ZF$
derivation.

The symmetric statement $\T \proves \varphi \Rightarrow \ZF \models \varphi$ needs
``every $\ZF$-model satisfies \cq{ClosureFO}'' --- the deep reverse direction,
blocked at the recursion theorem as discussed.

% =====================================================================
\section{Mechanization IV: G\"odel completeness, compactness, and the syntactic equivalence}\label{sec:mech-compl}
% =====================================================================

The fourth file, \texttt{Completeness.v}, proves the converse of soundness ---
\textbf{G\"odel completeness} --- from scratch, then leverages it to obtain
compactness for sentence theories and the forward \emph{syntactic} equivalence
$\ZF \proves \varphi \Rightarrow \T \proves \varphi$. (It builds on \texttt{Deep.v}
through Coq's namespacing: \cq{coqc -Q . SetTheory Deep.v} then
\cq{coqc -Q . SetTheory Completeness.v}.)

\subsection{The goal}

\begin{lstlisting}
Theorem completeness :
  forall G phi,
    (forall (Dom : Type) (m : Dom -> Dom -> Prop) (v : nat -> Dom),
       (forall g, In g G -> Sat Dom m v g) -> Sat Dom m v phi) ->
    Prov G phi.

Corollary prov_iff_valid :
  forall G phi,
    Prov G phi <->
    (forall (Dom : Type) (m : Dom -> Dom -> Prop) (v : nat -> Dom),
       (forall g, In g G -> Sat Dom m v g) -> Sat Dom m v phi).
\end{lstlisting}

In words: \emph{a formula is provable in the calculus if and only if it is valid in
every model.} The ``only if'' is soundness; the ``if'' is the content of
completeness. The standard route is Henkin's: show that a \emph{consistent} set of
formulas has a model (then contrapose). We outline the construction as it appears
in the file.

\subsection{Proof-theory infrastructure}

The file first develops the usual admissible rules: weakening, the deduction
theorem, proof by contradiction and double-negation elimination, a notion of
consistency \cq{Con G := ~ Prov G fBot}, the Lindenbaum extension step
\cq{Con_cons_or} (a consistent set stays consistent when extended by a formula or
its negation), the equality kit (symmetry/transitivity/congruence --- available
only \emph{because} the equality rule was corrected to the Leibniz \cq{P_eqElim}),
the admissibility of renaming \cq{Prov_rename}, and cut \cq{Prov_cut}.

\subsection{Model existence: the term model and the truth lemma}

The heart is \cq{model_exists}: an abstract \emph{maximal-consistent Henkin}
theory \cq{T} (one that decides every formula and provides a witness for every
existential) has a model. The model is the \emph{quotient term model}. Since terms
are just variables (\cq{nat}), the domain ought to be variables modulo provable
equality \cq{ceq i j}. To get an honest Coq type rather than a setoid, we pick
\emph{canonical representatives} with Hilbert's $\varepsilon$:

\begin{lstlisting}
Definition rep (i : nat) : nat := epsilon (inhabits 0) (fun j => ceq i j).
Definition D : Type := { n : nat | rep n = n }.       (* canonical reps *)
Definition memD (x y : D) : Prop := cmem (proj1_sig x) (proj1_sig y).
\end{lstlisting}

The linchpin is the \textbf{truth lemma}: in this model, satisfaction of a formula
coincides with its membership in the theory \cq{T}. It is proved by strong
induction on formula \emph{size} (\cq{fsize}), because the quantifier cases recurse
on an instantiated formula --- and instantiation is a renaming, which preserves
size:

\begin{lstlisting}
Lemma truth :
  forall n a, fsize a <= n ->
    forall s : nat -> nat,
      Sat D memD (fun i => mkD (s i)) a <-> T (rename s a).
\end{lstlisting}

The atomic cases use the congruence of \cq{cmem}/\cq{ceq} under the
representative map; the connective cases use that \cq{T} is maximal-consistent
(it respects $\to$, $\wedge$, $\vee$, $\neg$); the quantifier cases use that \cq{T}
is Henkin (an existential is in \cq{T} iff some witness instance is, and dually for
$\forall$).

\subsection{Lindenbaum and Henkin: building the maximal theory}

To apply \cq{model_exists} we must \emph{construct} a maximal-consistent Henkin
extension of any consistent set. This needs an enumeration of all formulas; we
build one from Cantor's pairing function and prove it surjective:

\begin{lstlisting}
Definition Enum (n : nat) : form := decode (S n) n.
Lemma Enum_surj : forall f, exists n, Enum n = f.
\end{lstlisting}

Then a chain \cq{chain G0 n} processes formulas one at a time: at each step it
decides the next formula (adding it or its negation, by \cq{Con_cons_or}) and, when
it adds an existential, also adds a \emph{fresh-witness} instance (the eigenvariable
and Henkin-witness lemmas, from \cq{Prov_rename} plus a freshness analysis of free
variables). The limit theory \cq{TL} is shown to satisfy all five hypotheses of
\cq{model_exists} --- consistent, complete, deductively closed, and Henkin for both
$\exists$ and $\forall$. Combining: \cq{model_of_con} (every consistent set has a
model), and contraposing gives \cq{completeness}, hence \cq{prov_iff_valid}.

\subsection{Compactness: infinite completeness for sentence theories}

Completeness as stated is about a \emph{finite} context \cq{G} (a list). To reach
genuinely infinite theories --- such as $\ZF$ and $\T$, whose Separation/Closure
schemas have infinitely many instances --- we need the compactness-style lift: if
every model of an infinite theory $B$ satisfies $\varphi$, then $B$ proves
$\varphi$ from a finite subset.

The obstacle is Henkin freshness. The witness-adding step needs a variable
\emph{fresh for the whole theory}; for an infinite theory that mentions all
variables, there is none. The resolution is to restrict to \emph{sentence}
theories --- theories all of whose axioms are sentences (no free variables):

\begin{lstlisting}
Definition Sentence (f : form) : Prop := forall n, ~ Free n f.
Definition Sentences (B : form -> Prop) : Prop := forall f, B f -> Sentence f.
Definition BProv (B : form -> Prop) (G : list form) (phi : form) : Prop :=
  exists Gb, (forall x, In x Gb -> B x) /\ Prov (Gb ++ G) phi.
\end{lstlisting}

A witness fresh for the \emph{finite} list of formulas added so far is then
automatically fresh with respect to the sentence base $B$ --- sentences contribute
no free variables to clash with. This yields:

\begin{lstlisting}
Theorem model_of_BCon :
  forall B L0, Sentences B -> BCon B L0 ->
    exists Dom m v, (forall g, B g -> Sat Dom m v g)
                 /\ (forall g, In g L0 -> Sat Dom m v g).

Theorem completeness_inf :
  forall B psi, Sentences B -> Sentence psi ->
    (forall Dom m v, (forall g, B g -> Sat Dom m v g) -> Sat Dom m v psi) ->
    BProv B nil psi.
\end{lstlisting}

and, as an immediate corollary, the model-theoretic characterization of deductive
equivalence:

\begin{lstlisting}
Theorem theory_equiv :
  forall B1 B2, Sentences B1 -> Sentences B2 ->
    (forall Dom m v, (forall g, B1 g -> Sat Dom m v g) <-> (forall g, B2 g -> Sat Dom m v g)) ->
    forall psi, Sentence psi -> (BProv B1 nil psi <-> BProv B2 nil psi).
\end{lstlisting}

\cq{theory_equiv} is the abstract engine behind the whole article:
\emph{two sentence theories with the same models prove the same sentences.}

\subsection{$\ZF$ and $\T$ as sentence theories, and $\ZF \proves \varphi \Rightarrow \T \proves \varphi$}

To feed $\ZF$ and $\T$ into this machinery we encode them as sentence theories.
Each axiom is universally closed by \cq{seal} (close every free variable with
$\forall$), so the parametrized schema instances become genuine sentences;
crucially, the Closure schema is encoded as a closed formula \cq{Closure_form}
whose \cq{Sat}-meaning is bridged back to the abstract \cq{SetLike}/closure
statement (\cq{bridge_Closure} and friends):

\begin{lstlisting}
Definition seal (f : form) : form := closeN (bound f) f.    (* universal closure *)
Definition Tax_s  (f : form) : Prop := (* sealed Ext, Reg, Pow, Sep-schema, Closure-schema *)
Definition ZFax_s (f : form) : Prop := (* sealed Ext, Reg, Pow, Pair, Union, Inf, Sep-schema, Repl-schema *)
\end{lstlisting}

The semantic content of the forward equivalence is repackaged as: \emph{every model
of $\T$ is a model of $\ZF$.}

\begin{lstlisting}
Lemma Tmodel_sat_ZF :
  forall (V : Type) (mem : V -> V -> Prop) v,
    (forall g, Tax_s g -> Sat V mem v g) -> (forall g, ZFax_s g -> Sat V mem v g).
\end{lstlisting}

The proof extracts the abstract axioms from a $\T$-model through the \cq{bridge_*}
lemmas, then reuses \texttt{Deep.v}'s derived satisfaction facts (\cq{sat_Pair},
\cq{sat_Union}, \cq{sat_Inf}, \cq{sat_Repl}) to discharge the four traded $\ZF$
axioms. Soundness plus \cq{completeness_inf} then deliver the forward
\emph{syntactic} direction:

\begin{lstlisting}
Theorem ZF_implies_T :
  forall phi, Sentence phi -> BProv ZFax_s nil phi -> BProv Tax_s nil phi.
\end{lstlisting}

Everything $\ZF$ proves, $\T$ proves. Read against \cq{theory_equiv}, this is one
half of the statement that $\ZF$ and $\T$ are the same theory --- and the half that
goes through cleanly.

% =====================================================================
\section{What is proven, and the one remaining direction}\label{sec:remaining}
% =====================================================================

\begin{center}
\renewcommand{\arraystretch}{1.3}
\begin{tabular}{@{}p{0.46\textwidth}p{0.30\textwidth}l@{}}
\toprule
Statement & Status & File \\
\midrule
$\T \models$ Pairing, Union, Replacement, Infinity & machine-checked & \texttt{Forward.v} \\
$\ZF \models$ Closure (set-like relations) & machine-checked & \texttt{Reverse.v} \\
forward trade from \emph{first-order} schemas & machine-checked & \texttt{Deep.v} \\
sound ND calculus; $\ZF\proves\varphi \Rightarrow \T\models\varphi$ & machine-checked & \texttt{Deep.v} \\
truth lemma / model existence & machine-checked & \texttt{Completeness.v} \\
completeness $\Prov G\,\varphi \Leftrightarrow G\models\varphi$ & machine-checked & \texttt{Completeness.v} \\
compactness ($B\models\varphi \Rightarrow B\proves\varphi$, sentences) & machine-checked & \texttt{Completeness.v} \\
same-model sentence theories are deductively equal & machine-checked & \texttt{Completeness.v} \\
$\ZF \proves \varphi \Rightarrow \T \proves \varphi$ & machine-checked & \texttt{Completeness.v} \\
$\T \proves \varphi \Rightarrow \ZF \proves \varphi$ & needs the recursion theorem & --- \\
\bottomrule
\end{tabular}
\end{center}

\subsection{The converse, characterized precisely}

The one direction we have \emph{not} formalized is the converse syntactic
statement $\T \proves \varphi \Rightarrow \ZF \proves \varphi$. It is true. By
\cq{theory_equiv} it would follow from ``every $\ZF$-model satisfies
\cq{ClosureFO},'' the dual of \cq{Tmodel_sat_ZF}. That is exactly the
\emph{deep reverse} direction --- and it is genuinely hard, for a reason that has
been visible since \S\ref{sec:mech-deep}.

\texttt{Reverse.v} proves Closure using a \emph{second-order} ingredient: the
iteration $W_n$ lives on the metatheory's \cq{nat}, and the family $\{W_n\}$ is
collected by an external recursion. A first-order $\ZF$-model is not handed such a
meta-recursion; it must \emph{build} the family internally, as a set-coded
function, using only first-order $\ZF$ derivations. Doing that is precisely the
content of the \textbf{recursion theorem} on $\omega$:

\begin{itemize}[leftmargin=1.4em]
\item encode functions as sets of Kuratowski pairs;
\item formalize the ``$n$-approximation'' of the iteration and prove existence and
  uniqueness of approximations by $\omega$-induction, \emph{inside} the object
  calculus \cq{Prov};
\item take the union of approximations to obtain the total iteration function, then
  its range.
\end{itemize}

None of this is conceptually mysterious --- it is standard textbook set theory ---
but rendering it as first-order \emph{object derivations} (rather than meta-level
Coq recursion) is a self-contained and heavy formalization in its own right. It is
the honest boundary of the present development: we prove the converse
\emph{semantically} (it is implicit in \texttt{Reverse.v}'s shallow model
argument) but not \emph{syntactically inside \cq{Prov}}, because the syntactic
content of the reverse direction simply \emph{is} the recursion theorem.

\subsection{Trust: axioms used, no admitted goals}

Running \cq{Print Assumptions} on the headline theorems reports only the standard
axioms of classical mathematics, and nothing else:

\begin{lstlisting}
classic                              (* excluded middle in Prop *)
constructive_indefinite_description  (* Hilbert epsilon / description *)
functional_extensionality
propositional_extensionality
proof_irrelevance
\end{lstlisting}

There is no \cq{Admitted} and no custom \cq{Axiom} anywhere in the four files.
These five are exactly the principles one expects when mechanizing classical set
theory with a quotient term model: classical logic, a description operator to turn
existence axioms into operations, and the extensionality/irrelevance principles
that make the quotient an honest type. Note that
\cq{constructive_indefinite_description} is \emph{metatheoretic} choice (Hilbert's
$\varepsilon$ in the ambient logic), not the set-theoretic Axiom of Choice --- which,
as \S\ref{sec:theories} explained, appears in neither $\T$ nor $\ZF$ here, so these
certificates do not list it. In particular, the metatheory is \emph{not} assumed to
be stronger than $\ZF$ --- the development does not, for instance, rely on a universe
of $\ZF$-sets supplied by a library.

% =====================================================================
\section{Conclusion}\label{sec:conclusion}
% =====================================================================

Two punchlines, one mathematical and one about formalization.

\paragraph{The mathematics.} Four of $\ZF$'s axioms --- Pairing, Union, Infinity,
Replacement --- are not four independent ideas but four instances of \emph{one}:
collect the predecessors of a seed under a set-like relation. The single Closure
schema packages that idea, and over the structural base it is interderivable with
all four. The non-obvious lesson is \emph{where the power comes from}: not from
Closure alone but from Closure resting on \emph{hosting} --- the one-line fact that
every set is a member of some set. That fact makes ``has a definable predecessor''
imply ``is set-like'' for free, which is what turns a statement about transitive
closures into general collection. And the proof assistant pins down exactly how much
is doing the work, by reading off which hypotheses each derivation discharged. The
verdict is bracing: of the ``structural'' axioms, \emph{none} is used in full.
Regularity and Choice are passengers, used in neither direction; Powerset is used
only through its trivial hosting shadow $\forall a\,\exists y\,(a\in y)$, and only
forward (the reverse needs no Powerset at all). The machine-checked
$\ZF - \text{Powerset} - \text{Regularity} \proves \text{Closure}$, and the
host-only forward derivations, make this precise: the genuine engine of the whole
equivalence is just $\{$Extensionality, Separation, Closure$\}$ together with
hosting. What looked at first like a Powerset/Regularity mirror was an artifact of a
convenient proof.

\paragraph{The formalization.} We did not stop at the equivalence. The development
ascends from a shallow model argument, to a deep embedding with genuinely
first-order schemas, to a sound and \emph{complete} proof calculus built from
scratch (truth lemma, Lindenbaum/Henkin, G\"odel completeness), to compactness for
sentence theories, to the abstract theorem that same-model sentence theories are
deductively equivalent --- of which ``$\ZF$ and $\T$ prove the same sentences'' is
an instance, with its forward half $\ZF \proves \varphi \Rightarrow \T \proves
\varphi$ now a machine-checked corollary. The single remaining direction is
characterized exactly: it is the recursion theorem rendered as a first-order object
derivation, the deep reverse obstacle, and nothing less. Everything else is done,
with no admitted goals and only the standard classical axioms --- a complete and
honest account of what it takes to certify that a slogan (``one closure schema'')
really is $\ZF$ in disguise.

\appendix
% =====================================================================
\section{File inventory and build instructions}
% =====================================================================

The development is four Rocq/Coq files under \texttt{src/SetTheory/}.

\begin{center}
\renewcommand{\arraystretch}{1.25}
\begin{tabular}{@{}lp{0.62\textwidth}@{}}
\toprule
File & Contents \\
\midrule
\texttt{Forward.v} & shallow embedding; $\{$Ext, Sep, Closure$\}$ + hosting derive Pairing/Union/Replacement/Infinity; linchpin \cq{host}/\cq{host_spec} (Powerset enters only via \cq{powerset_gives_hosting}); dependency-audit \cq{Check}s show the four free of Powerset. \\
\texttt{Reverse.v} & shallow embedding; $\ZF$ derives Closure via the transitive-closure recursion; numerals distinct Foundation-free (\cq{onat_trans}, \cq{onat_no_self}, \cq{onat_inj}); audit shows \emph{both} Powerset and Regularity unused. \\
\texttt{Deep.v} & deep embedding (\cq{form}, \cq{Sat}, \cq{rename}, \cq{Sat_rename}); first-order \cq{SeparationFO}/\cq{ClosureFO}; forward trade re-proved; ND calculus \cq{Prov}; \cq{soundness}; \cq{ZF_provable_holds_in_T}. \\
\texttt{Completeness.v} & from-scratch G\"odel \cq{completeness}/\cq{prov_iff_valid}; \cq{completeness_inf} (compactness); \cq{theory_equiv}; \cq{ZF_implies_T}. \\
\bottomrule
\end{tabular}
\end{center}

\texttt{Forward.v}, \texttt{Reverse.v}, and \texttt{Deep.v} are independent (no
inter-file \cq{Require}) and need only the standard library.
\texttt{Completeness.v} \cq{Require}s \cq{Deep} (hence the namespace flag) and
additionally uses the standard classical/extensionality modules.

\begin{lstlisting}[language={}]
# developed against Rocq 9.0.1
coqc Forward.v
coqc Reverse.v
coqc Deep.v
# Completeness.v builds on Deep via the SetTheory namespace:
coqc -Q . SetTheory Deep.v
coqc -Q . SetTheory Completeness.v
\end{lstlisting}

\end{document}
